\par\medskip
\noindent\textbf{Definición 5.5.}
Sea $X\xrightarrow{f}Y$ un morfismo esencialmente de presentación
finita. Si existe un entero fijo $n$ que satisface (1) o (2) de 5.4,
pongamos
\[
\underline I_X(X|Y)=\underline I_X^n(X|Y),
\qquad
J_X(X|Y)=J_X^n(X|Y).
\]

%%%%%%%%%%  scan idx 130  |  folio 119  %%%%%%%%%%

Esta noción se extiende naturalmente a los esquemas formales. Sea
$\mathfrak X$ un esquema formal noetheriano y $\mathcal E$ un
$\mathfrak X$-módulo coherente. Para todo abierto afín
$U\subset\mathfrak X$, definamos
\[
\underline I_{\mathfrak X}^{\sisi{P}{p}}(\mathcal E)(U)
=I_{\mathfrak X(U)}^{\sisi{P}{p}}\bigl(\mathcal E(U)\bigr).
\]
Si $\mathfrak X=\varinjlim_\nu X_\nu$, entonces
$\mathcal E=\varprojlim_\nu\mathcal E_\nu$, donde
\[
\mathcal E_\nu
=\mathcal E\otimes_{\mathcal O_{\mathfrak X}}\mathcal O_{X_\nu},
\]
y el carácter noetheriano da
\[
\underline I_{\mathfrak X}^{\sisi{P}{p}}(\mathcal E)
=\varprojlim_\nu\underline I_{X_\nu}^{\sisi{P}{p}}(\mathcal E_\nu).
\]

Sea ahora $S$ un esquema formal noetheriano y $\mathfrak X$ un
esquema formal $S$-ádico esencialmente de tipo finito. Sea
$\underline I_{\mathfrak X}$ el haz de ideales que define la
inmersión diagonal cerrada
\[
\Delta:\mathfrak X\longrightarrow
\mathfrak X\times_S\mathfrak X
\]
(donde el producto fibrado se toma en la categoría de los esquemas
formales), de modo que
\[
0\longrightarrow\underline I_{\mathfrak X}
\longrightarrow\mathcal O_{\mathfrak X\times_S\mathfrak X}
\longrightarrow\mathcal O_{\mathfrak X}\longrightarrow0
\]
es exacta. Definamos
\[
\Omega_{\mathfrak X|S}
=\underline I_{\mathfrak X}/\underline I_{\mathfrak X}^{2}.
\]
Escribiendo
\[
\mathfrak X=\varinjlim_\nu X_\nu,
\qquad X_\nu=\mathfrak X\times_SS_\nu,
\]
tenemos
\[
\Omega_{\mathfrak X|S}
\simeq\varprojlim_\nu\Omega_{X_\nu|S_\nu},
\]
por lo que $\Omega_{\mathfrak X|S}$ es un $\mathfrak X$-módulo
coherente. Para todo $S$-esquema formal $S'$ de tipo finito existe un
isomorfismo canónico
\[
\Omega_{\mathfrak X\times_SS'|S'}
\simeq
\mathcal O_{S'}\otimes_{\mathcal O_S}\Omega_{\mathfrak X|S}.
\]

\medskip
\noindent\textbf{Definición 5.6.}
Sea $S$ un esquema formal noetheriano y $\mathfrak X$ un esquema
formal $S$-ádico esencialmente de tipo finito. Definamos
\[
\underline I_{\mathfrak X}^p(\mathfrak X|S)
=\underline I_{\mathfrak X}^p(\Omega_{\mathfrak X|S}),
\]
y sea $J_{\mathfrak X}^p(\mathfrak X|S)$ el subesquema formal cerrado
de $\mathfrak X$ definido por este haz de ideales.

Como
$\Omega_{\mathfrak X|S}=\varprojlim_\nu\Omega_{X_\nu|S_\nu}$,
tenemos
\[
\underline I_{\mathfrak X}^p(\mathfrak X|S)
=\varprojlim_\nu\underline I_{X_\nu}^p(X_\nu|S_\nu)
\]
y
\[
J_{\mathfrak X}^p(\mathfrak X|S)
=\varinjlim_\nu J_{X_\nu}^p(X_\nu|S_\nu).
\]

%%%%%%%%%%  scan idx 131  |  folio 120  %%%%%%%%%%

El esquema formal $\mathfrak X_{(x)}$ es de nuevo $S$-ádico y
esencialmente de tipo finito, y existe un diagrama conmutativo
\[
\begin{tikzcd}[column sep=huge,row sep=large]
\mathfrak X_{(\widehat{x})}
  \arrow[rr,"i_{(\widehat{x})}"] \arrow[dr]
  && \mathfrak X \\
& \mathfrak X_{(x)} \arrow[ur,"i_{(x)}"'] & {}
\end{tikzcd}
\]
% Comprobación directa de la fuente a 1800 ppp: el subíndice superior
% izquierdo y la etiqueta de su flecha son ambos \widehat{x}; el subíndice
% inferior y su etiqueta son x sin sombrero.
Si $x$ es un punto aislado de $\mathfrak X$, entonces
$i_{(x)}$ e $i_{(\widehat{x})}$ son ambos isomorfismos sobre la
componente conexa de $\mathfrak X$ que contiene a $x$.

En particular, si
\[
|\mathfrak X|=\{x_1,\ldots,x_m\}
\]
es un espacio discreto finito, entonces existe una identificación
canónica
\[
\mathfrak X=\coprod_{1\leq i\leq m}\mathfrak X_{(x_i)}
\qquad\text{(unión disjunta)}.
\]
